% !TEX program = pdflatex
% !BIB program = bibtex
\documentclass[11pt]{article}

% Change "review" to "final" to generate the final (sometimes called camera-ready) version.
% Change to "preprint" to generate a non-anonymous version with page numbers.
\usepackage[review]{acl}

% Standard package includes
\usepackage{times}
\usepackage{latexsym}

% For proper rendering and hyphenation of words containing Latin characters (including in bib files)
\usepackage[T1]{fontenc}
% For Vietnamese characters
% \usepackage[T5]{fontenc}
% See https://www.latex-project.org/help/documentation/encguide.pdf for other character sets

% This assumes your files are encoded as UTF8
\usepackage[utf8]{inputenc}

% This is not strictly necessary, and may be commented out,
% but it will improve the layout of the manuscript,
% and will typically save some space.
\usepackage{microtype}

% This is also not strictly necessary, and may be commented out.
% However, it will improve the aesthetics of text in
% the typewriter font.
\usepackage{inconsolata}

%Including images in your LaTeX document requires adding
%additional package(s)
\usepackage{graphicx}
\usepackage{amsmath}
\usepackage{amssymb}
\usepackage{booktabs}
\usepackage{multirow}
\usepackage{float}
\usepackage{placeins}
\usepackage{stfloats}
\usepackage{adjustbox}
\usepackage[nameinlink,noabbrev]{cleveref}
\emergencystretch=2em
\hfuzz=20pt
\vfuzz=20pt
\hbadness=10000
\vbadness=10000
% If the title and author information does not fit in the area allocated, uncomment the following
%
%\setlength\titlebox{<dim>}
%
% and set <dim> to something 5cm or larger.

\usepackage{soul}
\usepackage{color, xcolor}
\definecolor{lightblue}{RGB}{0, 191, 255}
\sethlcolor{lightblue}
\soulregister\cite7
\soulregister\eqref7
\soulregister\ref7 
\newcommand{\mathcolorbox}[2]{\colorbox{#1}{$\displaystyle #2$}}
\newcommand{\clr}[1]{\textcolor{blue}{#1}}

\title{A Multi-cluster Boundary Learning Method for Out-of-Scope Intent Detection via MiniLM Embedding}

% \author{
%   \textbf{Yihong Xu\textsuperscript{1}},
%   \textbf{Mingyu Kang\textsuperscript{1*}},
%   \textbf{Linyuan L\"u\textsuperscript{1*}}\\
%   \textsuperscript{1}University of Science and Technology of China, Hefei, China\\
%   \texttt{xuyihong@mail.ustc.edu.cn}\\
%   \texttt{kangmingyu@ustc.edu.cn}\\
%   \texttt{linyuan.lv@ustc.edu.cn}
% }
% Author information can be set in various styles:
% For several authors from the same institution:
% \author{Author 1 \and ... \and Author n \\
%         Address line \\ ... \\ Address line}
% if the names do not fit well on one line use
%         Author 1 \\ {\bf Author 2} \\ ... \\ {\bf Author n} \\
% For authors from different institutions:
% \author{Author 1 \\ Address line \\  ... \\ Address line
%         \And  ... \And
%         Author n \\ Address line \\ ... \\ Address line}
% To start a separate ``row'' of authors use \AND, as in
% \author{Author 1 \\ Address line \\  ... \\ Address line
%         \AND
%         Author 2 \\ Address line \\ ... \\ Address line \And
%         Author 3 \\ Address line \\ ... \\ Address line}

% \author{First Author \\
%   Affiliation / Address line 1 \\
%   Affiliation / Address line 2 \\
%   Affiliation / Address line 3 \\
%   \texttt{email@domain} \\\And
%   Second Author \\
%   Affiliation / Address line 1 \\
%   Affiliation / Address line 2 \\
%   Affiliation / Address line 3 \\
%   \texttt{email@domain} \\}



%\author{
%  \textbf{First Author\textsuperscript{1}},
%  \textbf{Second Author\textsuperscript{1,2}},
%  \textbf{Third T. Author\textsuperscript{1}},
%  \textbf{Fourth Author\textsuperscript{1}},
%\\
%  \textbf{Fifth Author\textsuperscript{1,2}},
%  \textbf{Sixth Author\textsuperscript{1}},
%  \textbf{Seventh Author\textsuperscript{1}},
%  \textbf{Eighth Author \textsuperscript{1,2,3,4}},
%\\
%  \textbf{Ninth Author\textsuperscript{1}},
%  \textbf{Tenth Author\textsuperscript{1}},
%  \textbf{Eleventh E. Author\textsuperscript{1,2,3,4,5}},
%  \textbf{Twelfth Author\textsuperscript{1}},
%\\
%  \textbf{Thirteenth Author\textsuperscript{3}},
%  \textbf{Fourteenth F. Author\textsuperscript{2,4}},
%  \textbf{Fifteenth Author\textsuperscript{1}},
%  \textbf{Sixteenth Author\textsuperscript{1}},
%\\
%  \textbf{Seventeenth S. Author\textsuperscript{4,5}},
%  \textbf{Eighteenth Author\textsuperscript{3,4}},
%  \textbf{Nineteenth N. Author\textsuperscript{2,5}},
%  \textbf{Twentieth Author\textsuperscript{1}}
%\\
%\\
%  \textsuperscript{1}Affiliation 1,
%  \textsuperscript{2}Affiliation 2,
%  \textsuperscript{3}Affiliation 3,
%  \textsuperscript{4}Affiliation 4,
%  \textsuperscript{5}Affiliation 5
%\\
%  \small{
%    \textbf{Correspondence:} \href{mailto:email@domain}{email@domain}
%  }
%}

\begin{document}
\maketitle

\begin{abstract}
Intent detection is a critical task that bridges human intents and system actions in human-machine interaction systems. However, there still exist challenges for detecting out-of-scope (OOS) intents. (i) The traditional methods view the OOS intent detection as a multi-class classification, then the detection accuracy decreases as the class number of the known intents increases; (ii) LLM-embedding methods require large parameters, that makes them difficult to train and practically deploy. Thus, this work proposes a multi-cluster boundary learning method to detect OOS intents via MiniLM embedding (i.e., \texttt{all-MiniLM-L6-v2}) in an one-class classification workflow. The method learns the boundaries of multi-cluster embeddings generated by MiniLM from the training utterances, and then rejects the out-of-domain utterances as OOS intents. Experiments are conducted on public CLINC150, StackOverflow and Banking77 datasets. The results show that the method achieves the state-of-the-art OOS intent detection performance compared the other baselines. Ablation studies are also conducted and the results show that the used MiniLM can better adapt to the workflow and utterance embedding requirements. The code is available at supplementary materials.
\end{abstract}

\section{Introduction}
Intent detection is a critical yet challenging task for human-machine interaction systems. It maps user utterances to system actions, which bridges human intents and system behaviors~\citep{tur2010what,louvan2020recent,wolflein2025agents}. However, human intents are complex and cannot be fully enumerated in a closed-set intent detection system. That means if the system receives an out-of-scope (OOS) intent utterance but fails to reject it, that will trigger incorrect actions and induce harmful consequences, as shown in figure~\ref{fig:motivation}. Thus, a more important and difficult task is to detect the OOS intents and reject them in a timely manner, which is actually an open intent detection task~\citep{hoffman2024inferring, muzahid2024survey}.

\begin{figure}[t]
    \centering
    \includegraphics[width=\columnwidth]{figures/bg 02.png}
    \caption{A diagram of OOS intent detection}
    \label{fig:motivation}  
\end{figure}

\clr{There are three types of intent detection methods, including statistical methods, deep-learning methods and LLM-embedding methods. Among them, the statistical methods mainly use lexical or grammar-derived features with statistical classifiers, such as SVM~\citep{haffner2003optimizing} and task-dependent n-gram and Naive Bayes classifiers~\citep{wang2002combination}. But these methods rely heavily on predefined feature representations. Then, the deep-learning methods learn semantic representations directly from utterances, e.g., attention-based RNN~\citep{liu2016attention}, slot-gated networks~\citep{goo2018slot} and pretrained BERT~\citep{chen2019bert}. However, these classifiers are mainly trained under the closed-set setting and cannot explicitly reject OOS intents. Thus, several open intent detection methods introduce explicit rejection mechanisms, including class-wise open classification~\citep{shu2017doc}, distance-based scoring~\citep{lee2018simple}, contrastive representation learning~\citep{zeng2021modeling}, adaptive decision boundaries~\citep{zhang2021adaptive} and multi-granularity boundaries~\citep{li2025multi}. After that, LLM-embedding methods further use pretrained sentence embeddings and LLMs for semantic representation and OOS detection~\citep{zhang2024new,wang2024beyond}. Moreover, uncertainty-based routing combines lightweight classifiers with LLMs to reduce unnecessary LLM inference~\citep{arora2024intent,zaera2025efficient}. }

However, the existing methods still face two limitations. First, the OOS queries with similar semantic meaning as known intents are easily absorbed into the decision regions of the known intents~\citep{li2025multi}, but in fact they are OOS. Thus, if mix OOS labels into known intent labels, and simply view the OOS intent detection as a multi-class classification, the accuracy of OOS rejection would decrease as the class number increases. Second, although LLM-embedding methods improve the ability of semantic understanding, they still require very large parameters, that induces remarkable computational costs and high sensitivity of prompt designs~\citep{zaera2025efficient}. Thus, they are difficult to train and deploy in real-time resource-constrained dialogue systems.

To address these limitations, this work proposes a multi-cluster boundary learning method for OOS intent detection with a cascade workflow. The workflow decouples the complex open intent detection problem into multiple simpler stages. This allows the workflow to use a lightweight MiniLM instead of LLM wih large-scale parameters. Moreover, it reduces the OOS intent detection to a one-class classification, instead of a complex multi-class classification. Moreover, we found that MiniLM can embed the observation data into multiple clusters. Thus, if the boundaries of clusters are learned, the out-of-domain utterances can be identified as OOS intents.

Thus, the main contributions are as follows:
\begin{enumerate}
    \item[1.] A multi-cluster boundary learning method is proposed. This method learns the boundaries of multi-cluster embeddings generated by MiniLM from the training utterances, and then rejects the out-of-domain utterances as OOS intents. 
    
    \item[2.] A cascade workflow is proposed to first conduct OOS intent detection as an one-class classification, and then conduct known intent detection as a close-set multi-class classification. This means the OOS intent detection task can be addressed independently. 

    \item[3] Experiments are conducted on public real-world datasets for OOS intent detection task. The results show that the method achieves the state-of-the-art performance for OOS intent detection. Moreover, ablation studies are also conducted, and the results show that the used MiniLM, \texttt{all-MiniLM-L6-v2}, can better adapt to the workflow and utterance embedding requirements.  
\end{enumerate}


\section{Related Work new}
The open intent detection methods can be mainly categorized into three types, i.e., statistical methods, deep-learning methods and LLM-embedding methods, according to the problem settings and the types of classifiers.   

\subsection{Statistical Methods for Intent Detection}
The statistical methods mainly relied on lexical patterns and statistical classifiers to detect the intent classes. Among them, \citet{gorin1997how} identify the corresponding call type by learning salient lexical patterns from labeled spoken requests, but these methods hardly capture diverse expressions of the same intent. Then \citet{wang2002combination} propose a method combining rule-based grammars with statistical classifiers, but the grammar-based components remain domain-dependent and require domain-specific adaptation. To reduce the dependence, \citet{haffner2003optimizing} use salient grammar fragments as input features for SVMs and optimize the combination of binary classifiers for multi-class call classification. But these methods still rely on manually designed features. Thus, \citet{tur2010what} proposes an n-gram method to model local contextual dependencies, but feature sparsity limits its effectiveness on diverse intent expressions. Thus, \citet{sarikaya2011deep} introduce Deep Belief Networks to learn multi-layer representations from unlabeled data and fine-tune the model for natural language call routing, which moves intent detection from feature-based statistical models toward learned representations.

\subsection{Deep-Learning Methods for Intent Detection}
Then, deep-learning methods learn semantic representations directly from utterances instead of relying on predefined lexical features. E.g., \citet{liu2016attention} use RNN to capture the sequential dependencies for joint intent detection and slot filling, and \citet{goo2018slot} propose a slot-gated model that explicitly uses intent information to guide slot filling. Moreover, \citet{chen2019bert} introduce BERT for intent detection by using pretrained contextual representations. These methods improve the in-distribution accuracy, but the classification mode is still closed-set, and the classifier can only identify unknown utterances with known intent labels, and cannot reject the OOS intents~\citep{hendrycks2017baseline, guo2017calibration}. 

Thus, several methods introduce explicit rejection mechanisms for unknown inputs. \citet{hendrycks2017baseline} use the maximum softmax probability as a confidence score and reject inputs with low confidence. Moreover, \citet{lee2018simple} model class-conditional features with Gaussian distributions and use Mahalanobis distance to construct an OOD confidence score. However, these methods mainly derive the rejection decision from classifier scores and predefined feature-distribution assumptions. Then, \citet{shu2017doc} propose Deep Open Classification (DOC), which uses sigmoid classifiers and rejection thresholds for open intent detection. \citet{zeng2021modeling} further use supervised contrastive learning to reduce intra-class variance and enlarge inter-class variance. However, these methods hardly learn geometric acceptance regions for known intents. Thus, \citet{zhang2021adaptive} propose Adaptive Decision Boundary (ADB), which learns an adaptive spherical boundary for each known class. \citet{zhang2023learning} further introduce distance-aware representation learning before boundary optimization. But these methods still represent each known class with a single centroid and spherical boundary. To model more complex intent distributions, \citet{li2025multi} propose Multi-granularity Open Intent Classification via Adaptive Granular-Ball Decision Boundary (MOGB), which uses adaptive granular-ball clustering and multiple local centroids and radii to construct multi-granularity decision regions for each intent.

\subsection{LLM-Embedding Methods for Intent Detection}
Recent methods further use LLM-embedding models to learn semantic representations for intent detection and OOS rejection. \citet{zhang2024new} fine-tune sentence transformers for intent classification and introduce an embedding reconstruction loss to reduce the overlap between in-scope and OOS representations. However, the rejection performance still depends on the quality of the learned embedding space. Thus, recent studies explore large language models (LLMs) for stronger semantic reasoning. \citet{wang2024beyond} show that LLMs provide strong zero-shot and few-shot OOD detection, but remain less competitive than fully fine-tuned discriminative models under the full-resource setting. \citet{arora2024intent} further combine sentence transformers with LLMs through uncertainty-based routing to balance accuracy and inference latency. Similarly, \citet{zaera2025efficient} invoke a fine-tuned LLM only for samples with high classifier uncertainty. These methods reduce unnecessary LLM inference, but OOS rejection still relies on classifier uncertainty and additional LLM computation.

Then our previous work \citep{xu2026multi} proposes a lightweight multi-cluster boundary learning method based on MiniLM embeddings for OOS rejection. However, the MiniLM encoder remains fixed, so the representation space cannot be directly optimized for the OOS detection objective. 

Thus, the existing methods still face two problems. (i) the OOS intent detection is usually viewed as a multi-class classification and hardly separated from the known intent detection. Many existing methods perform OOS rejection and known-intent discrimination in the same learned representation space. (ii) The multi-agent LLMs require remarkably large computation cost. And due to the reason (i), the current routing mechanism is not accurate. 

\section{Method new}

\subsection{Cascade Workflow for Intent Detection}
\label{sec:cascade}

The proposed method uses a three-stage cascade: Gate $\rightarrow$ Router $\rightarrow$ Expert, as shown in Figure~\ref{fig:model_architecture}. Given an input utterance $x$, the Gate first determines whether $x$ belongs to the Known intent space or should be rejected as OOS. If $x$ is rejected, the Gate directly outputs an OOS label. Otherwise, $x$ is passed to the following modules. Then, the Router discriminates the coarse-grained domain of the remaining Known intents, and the utterance is delivered to the corresponding Expert for fine-grained intent classification. Since the Router and Experts only process accepted Known samples, their task is closed-set classification. Thus, the Router and Experts are parametrized by SmolLM-135M and supervisedly fine-tuned on Known-intent observations through low-rank adaptation (LoRA)~\citep{hu2022lora}.

\begin{figure*}[t]
    \centering
    \includegraphics[width=\textwidth]{figures/overview 01.eps}
    \caption{An overview of the cascaded workflow via MiniLM-based boundary learning.}
    \label{fig:model_architecture}
\end{figure*}

For an input utterance $x$, the Gate $G(\cdot)$ first determines whether $x$ belongs to the Known intent space:
\begin{equation}
z=G(x),\quad z\in\{\mathrm{Known},\mathrm{OOS}\},
\end{equation}
where $z$ is the binary decision. If $z=\mathrm{OOS}$, an OOS label is returned. Otherwise, the Router $R(\cdot)$ predicts the domain:
\begin{equation}
d=R(x),\quad d\in\mathcal{D},
\end{equation}
where $\mathcal{D}$ is the domain set. Then, the corresponding Expert $E_d(\cdot)$ classifies the intent:
\begin{equation}
y=E_d(x),\quad y\in\mathcal{Y}_d\subseteq\mathcal{Y}_{\mathrm{K}},
\end{equation}
where $\mathcal{Y}_{\mathrm{K}}$ is the Known intent label set and $\mathcal{Y}_d$ is the intent subset for domain $d$.

The proposed cascade separates OOS detection from Known-intent classification. The Gate only determines whether an input is supported by the Known intent space, while the Router and Experts further discriminate the accepted Known intents. Thus, OOS rejection does not need to be jointly optimized with the following fine-grained intent classifier, and a lightweight representation model can be used for the Gate stage.

\subsection{MiniLM-Based Intent Embedding}

Since the Gate only needs to distinguish OOS from Known samples, a lightweight MiniLM encoder, \texttt{all-MiniLM-L6-v2}, is used to construct the semantic representation space. The MiniLM maps each input utterance $x$ to an embedding with fixed dimension:
\begin{equation}
e_0=f_{\theta}(x),\quad e_0\in\mathbb{R}^{h},
\label{eq:embed}
\end{equation}
where $f_{\theta}(\cdot)$ denotes the MiniLM-based encoder parameterized by $\theta$, and $h$ denotes the embedding dimension.

A frozen MiniLM provides general sentence-level semantic representations, but its embedding space is not explicitly optimized for the boundary used in OOS rejection. Thus, we further adapt the MiniLM representation using only Known-intent samples. To retain the pretrained semantic information while adjusting the intent geometry, the embedding layer and the first four Transformer blocks are frozen, while the last two Transformer blocks are trainable. Moreover, a lightweight residual projection is introduced after mean pooling:
\begin{equation}
e=e_0+W_2\operatorname{GELU}\left(W_1\operatorname{LN}(e_0)\right),
\end{equation}
where $\operatorname{LN}(\cdot)$ denotes layer normalization, $W_1$ maps the 384-dimensional representation to a 256-dimensional hidden space, and $W_2$ maps it back to the original dimension. Then, the Gate applies $L_2$ normalization:
\begin{equation}
\bar{e}=\frac{e}{\|e\|_2}.
\end{equation}

The representation is optimized according to the geometric structure of Known intents. For each Known intent $y$, a class center $c_y$ is computed from its current normalized training embeddings. Let $\delta_{i,y}=\|\bar{e}_i-c_y\|_2^2$ denote the squared distance between the $i$-th sample and the center of intent $y$. The representation objective is defined as:
\begin{equation}
\mathcal{L}_{\mathrm{rep}}=\mathcal{L}_{\mathrm{ce}}+\alpha\mathcal{L}_{\mathrm{intra}}+\beta\mathcal{L}_{\mathrm{inter}},
\label{eq:rep_loss}
\end{equation}
where $\mathcal{L}_{\mathrm{ce}}$ is a center-based classification loss for discriminating Known intents, $\mathcal{L}_{\mathrm{intra}}$ reduces the distance between each sample and its own intent center, and $\mathcal{L}_{\mathrm{inter}}$ imposes a margin between the own-intent center and the nearest competing center. Specifically,
\begin{equation}
\mathcal{L}_{\mathrm{ce}}=-\frac{1}{N}\sum_i\log\frac{\exp(-\delta_{i,y_i}/T)}{\sum_{y\in\mathcal{Y}_{\mathrm{K}}}\exp(-\delta_{i,y}/T)},
\end{equation}
\begin{equation}
\mathcal{L}_{\mathrm{intra}}=\frac{1}{N}\sum_i\delta_{i,y_i},
\end{equation}
and
\begin{equation}
\mathcal{L}_{\mathrm{inter}}=\frac{1}{N}\sum_i\left[\delta_{i,y_i}-\min_{y\neq y_i}\delta_{i,y}+m\right]_{+},
\end{equation}
where $T$ is the temperature and $m$ is the inter-class margin. Thus, the adapted representation preserves the discrimination among Known intents while improving the geometric structure used for the following boundary learning. Note that, no real or pseudo OOS samples are introduced into the representation training.

The adapted embedding space provides the basis for boundary learning. However, utterances with the same intent may still contain different local structures due to their wording, expression and semantic focus. Thus, instead of restricting every intent to a single compact region, the formulation allows each Known intent to contain one or multiple local centroids. For intent $y$, its centroid set is defined as:
\begin{equation}
\mathcal{C}_{y}=\{c_{y,1},c_{y,2},\ldots,c_{y,K_y}\},
\label{eq:ky}
\end{equation}
where $c_{y,k},k=1,\dots,K_y$, denotes the $k$-th centroid of Known intent $y$, and $K_y$ is the number of centroids. When $K_y>1$, the local centroids are constructed by K-means clustering within each intent, while $K_y=1$ reduces to a single class centroid. Multiple centroids provide a more flexible representation of local structures, but the suitable number of centroids may vary with the embedding distribution. Thus, the influence of $K_y$ is further investigated in the experiments.

\subsection{Multi-cluster Boundary Learning for OOS Intent Detection}

The centroid representation defines local semantic regions for Known intents, and OOS detection is therefore formulated as a boundary learning problem in the embedding space. For each centroid, the Gate estimates a local boundary from its assigned Known samples. An input supported by at least one Known-intent region is accepted, while an input outside all learned regions is rejected as OOS.

For each centroid $c_{y,k}$, the Gate estimates a local radius $r_{y,k}$ from the Known training samples assigned to the corresponding region. One radius estimator is defined as:
\begin{equation}
r_{y,k}=\mu_{y,k}+\lambda\sigma_{y,k},
\label{eq:lambda}
\end{equation}
where $\mu_{y,k}$ and $\sigma_{y,k}$ are the mean and standard deviation of the within-region distances, and $\lambda$ controls the boundary width. Moreover, quantile-based radii are considered to reduce the dependence on a fixed distributional form. They estimate the local boundary from a selected quantile of the within-cluster or within-intent training distances.

The distance between an embedding $\bar{e}$ and a centroid $c_{y,k}$ can be measured by Euclidean distance or diagonal Mahalanobis distance. The latter is defined as:
\begin{equation}
d_{\mathrm{M}}(\bar{e},c_{y,k})=\sqrt{\sum_j\frac{(\bar{e}_j-c_{y,k,j})^2}{\sigma_{y,k,j}^{2}+\epsilon}},
\end{equation}
where $\sigma_{y,k,j}^{2}$ denotes the feature-wise variance estimated from the training samples assigned to the local region, and $\epsilon$ is a small regularization constant. The diagonal Mahalanobis distance accounts for feature-wise variation without estimating a full covariance matrix.

For an incoming utterance, the Gate first computes the normalized distance to each local region. Two acceptance rules are considered. The nearest-sphere rule first identifies the centroid with the minimum raw distance and normalizes its distance by the corresponding radius, while the normalized-union rule directly selects the minimum normalized distance:
\begin{equation}
s_{\mathrm{union}}(\bar{e})=\min_{y\in\mathcal{Y}_{\mathrm{K}},\,1\leq k\leq K_y}\frac{d(\bar{e},c_{y,k})}{r_{y,k}}.
\label{eq:ryk}
\end{equation}

Moreover, semantically ambiguous samples can be close to more than one Known intent. Thus, we additionally consider the relative distance between the nearest and second-nearest intent regions. Let $d_1$ and $d_2$ denote the nearest and second-nearest distances, respectively. Their ratio $d_1/d_2$ and inverse margin $d_1/(d_2-d_1)$ are used as optional margin terms and can be combined with the normalized boundary score. The candidate fusion rules include the normalized score alone, additive ratio fusion, additive inverse-margin fusion and max-ratio fusion.

The final Gate decision is:
\begin{equation}
G(x)=
\begin{cases}
\mathrm{Known}, & s(\bar{e})\leq\tau,\\
\mathrm{OOS}, & s(\bar{e})>\tau,
\end{cases}
\end{equation}
where $s(\bar{e})$ denotes the selected boundary score and $\tau$ denotes the corresponding rejection threshold. The centroid number, distance metric, radius estimator, acceptance rule, fusion function and threshold are selected using Known validation samples. Thus, real or pseudo OOS samples are not required for boundary selection.

For $M$ Known intents, $K$ centroids per intent and an embedding dimension of $h$, the basic boundary search requires $O(MKh)$ distance operations for each input, while storing the centroids and their diagonal statistics requires $O(MKh)$ additional space.


\section{Experiments new}
\subsection{Datasets}

Three datasets are used for performance evaluation, as shown in Table~\ref{tab:dataset_stats}. The datasets are divided into training, validation and test sets according to their corresponding experimental protocols, and different Known Intent Ratios (KIRs) are used to construct the open intent detection settings.

\textbf{CLINC150} is a large-scale multi-domain dataset with 150 in-scope intents and explicit OOS samples~\citep{larson2019evaluation}. The data source is available at \url{https://github.com/clinc/oos-eval}.

\textbf{StackOverflow} is a technical-domain short-text dataset with 20 intent classes and substantial semantic overlap across classes. The data source is available at \url{https://github.com/jacoxu/StackOverflow}.

\textbf{Banking77} is a fine-grained, single-domain dataset with 77 intent classes. It focuses on high-granularity classification in the financial domain. The data source is available at \url{https://github.com/PolyAI-LDN/task-specific-datasets}.

\begin{table}[H]
\centering
\renewcommand{\arraystretch}{1.05}
\resizebox{\columnwidth}{!}{
\begin{tabular}{lccc}
\toprule
\textbf{Dataset} & \textbf{\#Intents} & \textbf{\#Samples} & \textbf{Type} \\
\midrule
CLINC150 & 150 & 22500 & Multiple domains \\
StackOverflow & 20 & 20000 & Technical \\
Banking77 & 77 & 13083 & Banking \\
\bottomrule
\end{tabular}
}
\caption{Dataset statistics.}
\label{tab:dataset_stats}
\end{table}

\subsection{Baselines and Experimental Settings}

This work defines Known Intent Ratio (KIR) to describe the proportion of all intent classes that are treated as Known intents during training, as follows:
\begin{equation}
\mathrm{KIR}=\frac{|\mathcal{Y}_{\mathrm{K}}|}{|\mathcal{Y}|},
\end{equation}
where $\mathcal{Y}$ is the full intent label set and $\mathcal{Y}_{\mathrm{K}}$ is the subset of intents treated as Known during training.

Moreover, the following representative methods are selected as baselines:

\textbf{MSP}~\citep{hendrycks2017baseline} uses the maximum softmax probability as the confidence score for detecting misclassified or out-of-distribution samples.

\textbf{OpenMax}~\citep{bendale2016towards} recalibrates activation scores and estimates the probability that an input belongs to an unknown class.

\textbf{DOC}~\citep{shu2017doc} replaces the softmax layer with independent sigmoid classifiers and performs class-wise open-set rejection.

\textbf{DeepUnk}~\citep{lin2019deep} uses margin loss to learn discriminative intent features and applies novelty detection for unknown intent detection.

\textbf{KNNCL}~\citep{zhou2022knn} combines KNN-based contrastive representation learning with nearest-neighbor information for OOD intent detection.

\textbf{ADB}~\citep{zhang2021adaptive} learns adaptive class-specific decision boundaries for open intent detection.

\textbf{DA-ADB}~\citep{zhang2023learning} improves ADB with distance-aware representation learning and adaptive boundary learning.

\textbf{MOGB}~\citep{li2025multi} models each Known intent with adaptive granular-ball representations and constructs multiple local decision boundaries. Since MOGB also considers local structures within each intent, we further compare it with the proposed boundary learning method in the ablation experiments.

Moreover, the Gate uses the 22M-parameter \texttt{all-MiniLM-L6-v2}, while the Router and Expert use \texttt{SmolLM-135M}. For the Gate, only the last two Transformer blocks and the residual projection are trainable, with learning rates of $2\times10^{-5}$ and $2\times10^{-4}$, respectively. The intra-class and inter-class weights in Eq.~(\ref{eq:rep_loss}) are both set to 0.1, with temperature 0.07 and margin 0.20. For the Router and Expert, the LoRA ranks are 32 and 16, respectively, and the learning rate is $2\times10^{-4}$. All main experiments use only Known training and validation samples for Gate training and model selection, while test samples are used only for final evaluation. The OOS-aware selection variant is used only in the ablation study. All main results are averaged over three random seeds, i.e., 13, 42 and 87.


\subsection{Evaluation Metrics}

To evaluate and compare the performances of all methods, Known F1 score, OOS F1 score and Acc are selected as metrics.

\textbf{Known F1 score}~\citep{wang2021texsmart,zawbaa2024improved} is calculated by averaging F1 scores over all Known intent classes, as follows:
\begin{equation}
\mathrm{Known\ F1}=\frac{1}{|\mathcal{Y}_{\mathrm{K}}|}\sum_{y\in\mathcal{Y}_{\mathrm{K}}}\frac{2P_yR_y}{P_y+R_y},
\end{equation}
where $\mathcal{Y}_{\mathrm{K}}$ denotes the Known-intent label set, and $P_y$ and $R_y$ denote the precision and recall of class $y$, respectively. The precision is calculated from the complete open-set test predictions, so OOS samples incorrectly accepted as Known intents are also counted as false positives of the corresponding Known classes.

\textbf{OOS F1 score}~\citep{wang2021texsmart,zawbaa2024improved} is calculated for evaluating OOS detection performance, as follows:
\begin{equation}
\mathrm{OOS\ F1}=\frac{2P_{\mathrm{OOS}}R_{\mathrm{OOS}}}{P_{\mathrm{OOS}}+R_{\mathrm{OOS}}},
\end{equation}
where
\begin{equation}
\begin{aligned}
P_{\mathrm{OOS}}&=\frac{\mathrm{TP}_{\mathrm{OOS}}}{\mathrm{TP}_{\mathrm{OOS}}+\mathrm{FP}_{\mathrm{OOS}}},\\
R_{\mathrm{OOS}}&=\frac{\mathrm{TP}_{\mathrm{OOS}}}{\mathrm{TP}_{\mathrm{OOS}}+\mathrm{FN}_{\mathrm{OOS}}}.
\end{aligned}
\end{equation}
$\mathrm{TP}_{\mathrm{OOS}}$, $\mathrm{FP}_{\mathrm{OOS}}$ and $\mathrm{FN}_{\mathrm{OOS}}$ denote true positives, false positives and false negatives w.r.t. the OOS class, respectively.

\textbf{Acc} is the overall accuracy over all test samples, as follows:
\begin{equation}
\mathrm{Acc}=\frac{1}{N}\sum_{i=1}^{N}\mathbb{I}(\hat{y}_i=y_i),
\end{equation}
where $N$ is the size of the test set, $y_i$ and $\hat{y}_i$ are the ground truth and predicted labels of the $i$-th sample, respectively, and $\mathbb{I}(\cdot)$ is the indicator function.


\subsection{Performance on OOS Intent Detection}

The performance on OOS intent detection is presented in Table~\ref{tab:main_results_all}. For different datasets and KIR settings, the proposed method selects the boundary configuration using Known validation samples while retaining the same Trainable MiniLM Gate and Gate--Router--Expert cascade workflow. The results of our method are reported over three random seeds.

\begin{table*}[!t]
\centering
\tiny
\setlength{\tabcolsep}{2.8pt}
\renewcommand{\arraystretch}{0.97}
\resizebox{\textwidth}{!}{
\begin{tabular}{c|l|ccc|ccc|ccc}
\toprule
\multirow{2}{*}{\textbf{KIR}} & \multirow{2}{*}{\textbf{Method}} & \multicolumn{3}{c|}{\textbf{CLINC150}} & \multicolumn{3}{c|}{\textbf{StackOverflow}} & \multicolumn{3}{c}{\textbf{Banking77}}\\
\cmidrule{3-11}
& & \textbf{Known F1} & \textbf{OOS F1} & \textbf{Acc} & \textbf{Known F1} & \textbf{OOS F1} & \textbf{Acc} & \textbf{Known F1} & \textbf{OOS F1} & \textbf{Acc}\\
\midrule
\multirow{8}{*}{$0.25$}
& MSP & 51.02 & 59.26 & 53.38 & 42.66 & 11.66 & 27.91 & 50.47 & 39.42 & 42.19\\
& OpenMax & 62.65 & 77.51 & 70.27 & 47.51 & 34.52 & 38.97 & 53.42 & 48.52 & 47.76\\
& DOC & 75.46 & 90.78 & 86.08 & 56.30 & 62.50 & 57.75 & 65.16 & 76.64 & 70.31\\
& DeepUnk & 76.95 & 91.61 & 87.18 & 47.39 & 36.87 & 40.03 & 64.97 & 76.98 & 70.68\\
& KNNCL & 78.85 & 93.56 & 89.87 & 41.79 & 15.26 & 28.65 & 65.54 & 79.34 & 73.01\\
& ADB & 77.85 & 92.36 & 88.30 & 77.62 & 90.96 & 86.75 & 70.92 & 85.05 & 79.33\\
& DA-ADB & \textbf{79.57} & 93.20 & 89.48 & 80.87 & 92.65 & 89.07 & \textbf{73.05} & 86.57 & 81.19\\
& \textbf{Ours} & 75.39$\pm$1.65 & \textbf{95.19$\pm$0.48} & \textbf{91.19$\pm$0.27} & \textbf{82.61$\pm$0.13} & \textbf{95.50$\pm$0.96} & \textbf{91.58$\pm$0.92} & 62.83$\pm$2.41 & \textbf{92.58$\pm$0.89} & \textbf{86.91$\pm$1.36}\\
\midrule
\multirow{8}{*}{$0.50$}
& MSP & 72.82 & 63.71 & 66.68 & 66.28 & 26.94 & 53.23 & 73.20 & 46.29 & 61.67\\
& OpenMax & 79.83 & 82.15 & 80.22 & 69.88 & 46.11 & 60.27 & 75.16 & 55.03 & 65.53\\
& DOC & 83.84 & 87.45 & 85.19 & 77.37 & 71.18 & 73.88 & 78.38 & 72.66 & 74.60\\
& DeepUnk & 83.30 & 87.48 & 84.95 & 67.67 & 35.80 & 55.46 & 75.61 & 67.80 & 71.01\\
& KNNCL & 83.25 & 87.85 & 85.32 & 61.50 & 8.50 & 45.38 & 75.16 & 67.21 & 70.41\\
& ADB & 85.12 & 88.60 & 86.54 & 85.32 & 87.70 & 86.49 & 81.39 & 79.43 & 79.61\\
& DA-ADB & \textbf{85.58} & 90.10 & \textbf{87.93} & \textbf{86.71} & 88.86 & \textbf{87.78} & \textbf{82.54} & 79.93 & \textbf{81.51}\\
& \textbf{Ours} & 76.98$\pm$0.92 & \textbf{92.13$\pm$0.81} & 84.96$\pm$0.81 & 83.38$\pm$1.60 & \textbf{89.24$\pm$2.03} & 85.89$\pm$1.80 & 65.86$\pm$1.29 & \textbf{88.47$\pm$0.67} & 81.25$\pm$0.72\\
\midrule
\multirow{8}{*}{$0.75$}
& MSP & 83.65 & 63.86 & 76.19 & 81.42 & 37.86 & 73.20 & 84.99 & 46.05 & 77.08\\
& OpenMax & 71.14 & 75.18 & 75.36 & 82.98 & 49.69 & 75.78 & 85.50 & 53.02 & 78.32\\
& DOC & 87.91 & 83.87 & 85.93 & 85.64 & 65.32 & 80.55 & 84.14 & 63.51 & 78.94\\
& DeepUnk & 86.57 & 82.67 & 84.61 & 80.51 & 34.38 & 71.56 & 81.65 & 50.57 & 74.73\\
& KNNCL & 86.14 & 82.05 & 84.12 & 76.16 & 7.19 & 65.01 & 81.76 & 51.42 & 74.78\\
& ADB & \textbf{88.97} & 84.85 & 86.99 & 86.91 & 74.10 & 82.89 & \textbf{86.44} & 67.34 & \textbf{81.39}\\
& DA-ADB & 88.43 & 86.00 & \textbf{87.39} & \textbf{87.66} & 74.55 & \textbf{83.56} & 85.93 & 69.37 & 81.12\\
& \textbf{Ours} & 78.05$\pm$0.38 & \textbf{86.80$\pm$1.55} & 79.56$\pm$0.51 & 84.24$\pm$0.91 & \textbf{75.95$\pm$1.94} & 81.92$\pm$0.96 & 70.09$\pm$0.25 & \textbf{86.60$\pm$0.42} & 79.44$\pm$0.38\\
\bottomrule
\end{tabular}
}
\caption{Performance comparison under different KIRs. All OOS F1 results of our method follow Known-only training and selection and are reported over three random seeds.}
\label{tab:main_results_all}
\end{table*}

As shown in Table~\ref{tab:main_results_all}, the OOS F1 scores of most methods decrease as KIR increases, indicating that OOS rejection becomes more difficult when more intent classes are included in the Known intent space. But the proposed method maintains competitive OOS detection performance under different KIRs and achieves the highest OOS F1 among the compared methods in all nine settings.

Moreover, the proposed method does not always achieve the highest Known F1. This indicates that OOS rejection and fine-grained Known-intent classification have different optimization objectives. In the proposed cascade, the Gate is responsible for detecting whether an utterance is supported by the Known-intent space, while the Router and Experts further classify the accepted samples. Thus, the downstream classifier can be replaced by a stronger Known-intent classification method without changing the OOS detection mechanism.

\subsection{Interpretability Analysis for Gate Stage}

The quantitative results show the final OOS rejection performance, but they do not directly explain how representation adaptation changes the Gate decision. Thus, we first visualize a representative local boundary on StackOverflow when KIR$=0.50$, as shown in Figure~\ref{fig:local_boundary_geometry}.

\begin{figure*}[t]
    \centering
    \includegraphics[width=0.88\textwidth]{figures/paper_style_local_boundary_geometry_3d.png}
    \caption{Local boundary visualization on StackOverflow when KIR$=0.50$. The translucent surface denotes the normalized boundary $s(\bar{e})=1$. The three-dimensional coordinates are used only for visualization, while the actual Known/OOS decisions are obtained from the original 384-dimensional embeddings.}
    \label{fig:local_boundary_geometry}
\end{figure*}

As shown in Figure~\ref{fig:local_boundary_geometry}, some OOS samples are close to the Known-intent support region under the frozen MiniLM representation. After representation adaptation, the relative positions between these samples and the corresponding Known-intent boundary change, and more OOS samples fall outside the acceptance region. Note that, the dimensionality reduction is only used to display the local geometry and does not participate in the actual Gate decision.

A more direct observation can be obtained from the normalized Gate score, as shown in Figure~\ref{fig:gate_score_distribution}. Both Frozen and Trainable MiniLM use the same rejection threshold $s(\bar{e})=1$. Thus, the score distribution reflects whether representation adaptation changes the positions of OOS samples relative to the same decision boundary.

\begin{figure}[t]
    \centering
    \includegraphics[width=\columnwidth]{figures/paper_style_score_distribution.png}
    \caption{Gate score distributions of Frozen and Trainable MiniLM on StackOverflow when KIR$=0.50$. The dashed line denotes the fixed normalized threshold $s(\bar{e})=1$.}
    \label{fig:gate_score_distribution}
\end{figure}

As shown in Figure~\ref{fig:gate_score_distribution}, the Trainable MiniLM moves more OOS samples to the rejection side of the fixed boundary. Thus, the improvement does not simply result from changing the rejection threshold. Instead, the Known-only representation adaptation changes the normalized distances between the utterance embeddings and the learned Known-intent regions. But some overlap still remains around the threshold. Thus, OOS samples close to Known-intent boundaries remain a difficult case for the proposed method.

The error analysis further shows that near-OOS samples are substantially more difficult than OOS samples that are far from the Known-intent space. This pattern appears across all three datasets, while StackOverflow shows a particularly weak local separation between semantically adjacent intents. This is consistent with its larger performance degradation when KIR increases. Therefore, the remaining errors mainly arise from semantically adjacent Known and OOS samples rather than from clearly separated OOS samples.

\subsection{MiniLM Selection for Gate Stage}

We first examine whether the explicit Gate is necessary for the cascade workflow. Table~\ref{tab:ablation_all_kir} follows the original cascade ablation and compares the complete cascade with three variants. \textbf{Without Gate} removes the explicit geometric Gate and performs OOS rejection using the downstream Expert confidence. Specifically, the rejection score is defined as one minus the maximum Expert confidence, and its threshold is selected on the validation set. \textbf{Cascade-MiniLM} uses MiniLM for the cascade modules, while \textbf{Cascade-SmolLM} uses SmolLM for the Gate and downstream modules.

\begin{table*}[t]
\centering
\small
\setlength{\tabcolsep}{0pt}
\renewcommand{\arraystretch}{1.0}
\begin{tabular*}{\textwidth}{@{\extracolsep{\fill}}clcccccc@{}}
\toprule
\multirow{2}{*}{\textbf{KIR}} & \multirow{2}{*}{\textbf{Variant}} & \multicolumn{2}{c}{\textbf{CLINC150}} & \multicolumn{2}{c}{\textbf{StackOverflow}} & \multicolumn{2}{c}{\textbf{Banking77}}\\
\cmidrule(lr){3-4}
\cmidrule(lr){5-6}
\cmidrule(lr){7-8}
& & \textbf{Acc} & \textbf{OOS F1} & \textbf{Acc} & \textbf{OOS F1} & \textbf{Acc} & \textbf{OOS F1}\\
\midrule
\multirow{5}{*}{$0.25$}
& Ours & 91.19 & 95.19 & 91.58 & 95.50 & 86.91 & 92.58\\
& Frozen MiniLM & 89.83 & 94.29 & 89.30 & 93.85 & 85.33 & 91.58\\
& Without Gate & 73.00 & 81.75 & 66.84 & 74.95 & 82.72 & 89.81\\
& Cascade-MiniLM & 92.05 & 95.21 & 92.74 & 95.99 & 93.58 & 96.57\\
& Cascade-SmolLM & 80.85 & 88.09 & 40.33 & 45.10 & 51.15 & 64.41\\
\midrule
\multirow{5}{*}{$0.50$}
& Ours & 84.96 & 92.13 & 85.89 & 89.24 & 81.25 & 88.47\\
& Frozen MiniLM & 83.08 & 88.02 & 76.74 & 79.02 & 76.75 & 84.82\\
& Without Gate & 73.00 & 76.03 & 78.36 & 80.16 & 82.43 & 88.75\\
& Cascade-MiniLM & 89.67 & 91.64 & 89.80 & 91.85 & 85.93 & 91.46\\
& Cascade-SmolLM & 67.65 & 69.71 & 44.41 & 6.38 & 47.89 & 55.38\\
\midrule
\multirow{5}{*}{$0.75$}
& Ours & 79.56 & 86.80 & 81.92 & 75.95 & 79.44 & 86.60\\
& Frozen MiniLM & 77.76 & 78.95 & 75.17 & 63.13 & 76.16 & 83.55\\
& Without Gate & 71.65 & 61.76 & 77.40 & 65.94 & 76.00 & 82.32\\
& Cascade-MiniLM & 81.89 & 81.33 & 84.76 & 78.45 & 82.70 & 88.49\\
& Cascade-SmolLM & 68.44 & 55.02 & 64.86 & 0.13 & 45.88 & 42.07\\
\bottomrule
\end{tabular*}
\caption{Ablation results of the cascade workflow under different KIR settings. }
\label{tab:ablation_all_kir}
\end{table*}

As shown in Table~\ref{tab:ablation_all_kir}, removing the Gate causes a clear degradation in OOS detection, showing that the explicit rejection stage cannot be replaced simply by the confidence of the downstream classifier. Moreover, Cascade-SmolLM performs substantially worse in most settings, indicating that a larger language model does not necessarily provide a better representation for the boundary-based Gate. In contrast, Cascade-MiniLM obtains similar OOS F1 to Ours in several settings because they share the same MiniLM-based rejection mechanism, while their Acc differs due to the different downstream classifiers. This further shows that the Gate mainly determines OOS rejection, whereas the Router and Experts mainly affect the classification of accepted Known samples.

The Frozen MiniLM comparison further examines the effect of representation adaptation. The degradation is mainly associated with increased OOS false acceptance rather than a substantial improvement in Known coverage. This effect is particularly evident on StackOverflow at larger KIRs, where the frozen representation admits considerably more OOS samples into the Known regions. Together with the matched $K=1$ Frozen-versus-Trainable analysis, this indicates that adapting MiniLM mainly changes the representation geometry used by the boundary rather than simply shifting the rejection threshold.

Then, we investigate whether increasing the number of centroids consistently improves the boundary. Table~\ref{tab:k_ablation} reports the Gate-level OOS F1 when $K$ varies from 1 to 5 under KIR$=0.50$.

\begin{table*}[t]
\centering
\small
\setlength{\tabcolsep}{4.0pt}
\renewcommand{\arraystretch}{1.05}
\begin{tabular}{llccccc}
\toprule
\textbf{Dataset} & \textbf{Distance} & $K=1$ & $K=2$ & $K=3$ & $K=4$ & $K=5$\\
\midrule
CLINC150 & Euclidean & 87.70$\pm$1.33 & \textbf{88.42$\pm$1.27} & 88.18$\pm$0.69 & 87.45$\pm$0.58 & 87.08$\pm$0.37\\
CLINC150 & Diag. Mahalanobis & 88.02$\pm$1.00 & \textbf{88.07$\pm$1.06} & 87.33$\pm$0.61 & 86.12$\pm$0.61 & 85.39$\pm$0.46\\
Banking77 & Euclidean & 80.06$\pm$1.40 & 82.26$\pm$1.83 & 83.58$\pm$0.90 & 84.54$\pm$0.55 & \textbf{85.60$\pm$0.59}\\
Banking77 & Diag. Mahalanobis & 84.82$\pm$1.44 & 85.57$\pm$1.37 & 87.72$\pm$0.48 & 89.25$\pm$0.57 & \textbf{89.78$\pm$0.64}\\
StackOverflow & Euclidean & \textbf{77.36$\pm$5.91} & 64.80$\pm$7.49 & 55.19$\pm$2.70 & 52.90$\pm$5.54 & 55.67$\pm$4.88\\
StackOverflow & Diag. Mahalanobis & \textbf{79.02$\pm$4.63} & 72.80$\pm$7.02 & 66.21$\pm$5.27 & 65.02$\pm$7.58 & 66.80$\pm$5.46\\
\bottomrule
\end{tabular}
\caption{OOS F1 (\%) with different numbers of centroids under KIR$=0.50$.}
\label{tab:k_ablation}
\end{table*}

As shown in Table~\ref{tab:k_ablation}, the influence of $K$ is strongly dataset dependent. On CLINC150, increasing $K$ provides only a limited improvement before the performance decreases, while StackOverflow consistently degrades when more local centers are introduced. In contrast, Banking77 benefits from additional centers under this setting. This difference indicates that increasing $K$ changes not only the representation granularity but also the total acceptance region of the Gate. When the additional local regions mainly improve the coverage of heterogeneous Known samples, multiple centers can be beneficial; but when they expand toward nearby OOS samples, false acceptance increases instead. Thus, a fixed $K$ is not suitable for all datasets, which motivates selecting the boundary configuration according to the validation distribution rather than assuming that more centers are always better.

Moreover, MOGB~\citep{li2025multi} also represents Known intents with multiple local regions. To reduce the influence of different encoders, we further compare the two boundary designs under the same MiniLM representation at KIR$=0.50$. In this controlled comparison, the Trainable $K=1$ boundary obtains OOS F1 scores of 90.44\%, 87.67\% and 83.56\% on CLINC150, StackOverflow and Banking77, respectively, while the corresponding MOGB-MiniLM component obtains 81.32\%, 72.92\% and 74.99\%. Thus, the performance difference cannot be explained only by whether multiple local centers are used. But this comparison controls the representation backbone and boundary component rather than reproducing the complete original MOGB system, so it is not used for an unconditional superiority claim.

Finally, the Trainable MiniLM Gate contains about 22.9M parameters with approximately 3.75M trainable parameters. On the NVIDIA RTX 5070 used in our experiments, the measured batch-1 Gate latency is approximately 3.02 ms, 3.07 ms and 2.68 ms on CLINC150, StackOverflow and Banking77, respectively. Together with the $O(MKh)$ boundary complexity, these results characterize the additional computational cost introduced by the Gate. However, the measured latency is device-specific and is not used to claim universal deployment superiority.


\section{Conclusion}

This work proposes a multi-cluster boundary learning method for OOS intent detection via a cascade MiniLM workflow. That workflow separates the OOS intent detection from the multi-class classification of known intents, and views it as a one-class classification problem. Then, multi-cluster boundary learning is conducted to learn the boundary of MiniLM embeddings from the input user utterances. That boundary separates the embedding domain of OOS intents and known intents. Extensive experiments are conducted on public datasets. The results show that the proposed method achieve the stete-of-the-art OOS intent detection performance. Ablation studies also show that the MiniLM settings are reasonable at present. Moreover, the performance of the method is not always the best one on the task of known intent detection, but the cascade workflow completely decouples the two detection process for OOS and known intents. Thus, the result does not influence the aforementioned conclusion. And moreover, the lightweight models have outstanding advantages for real-world applications compared to the large models. The more advanced MiniLMs perhaps improve the detection performance by using this framework in the future.       




\section*{Limitations}
The proposed method in this work achieves the state-of-the-art performance on OOS intent detection task, but does not generalize well on the task of known intent detection. Thus our primary problem setting is OOS intent detection in the title. 

\section*{Ethics Statement}
This paper does not involve text reuse from previously published work. We do not use any AI-assisted tools for writing this paper.

\bibliography{custom}

\end{document}
